% !TEX program = lualatex
% Prueba 1 de LuaLaTeX en TeXsting: Lua haciendo cálculos dentro del documento.
\documentclass[11pt,a4paper]{article}
\usepackage[spanish]{babel}
\usepackage{fontspec}
\usepackage{unicode-math}
\setmainfont{Latin Modern Roman}
\setmathfont{Latin Modern Math}
\usepackage{luacode}
\usepackage{booktabs}
\usepackage[margin=2.2cm]{geometry}

\begin{luacode*}
  -- Funciones en Lua que el documento usa más abajo
  function primos(n)
    local lista, es = {}, {}
    for i = 2, n do es[i] = true end
    for i = 2, n do
      if es[i] then
        lista[#lista + 1] = i
        for j = i * i, n, i do es[j] = false end
      end
    end
    tex.sprint(table.concat(lista, ", "))
  end

  function fibonacci(n)
    local a, b, t = 0, 1, {}
    for i = 1, n do t[#t + 1] = a; a, b = b, a + b end
    tex.sprint(table.concat(t, ", "))
  end

  function tabla_trig()
    for g = 0, 90, 15 do
      local r = math.rad(g)
      local tg = (g == 90) and "---" or string.format("%.4f", math.tan(r))
      tex.sprint(string.format("%d° & %.4f & %.4f & %s \\\\", g, math.sin(r), math.cos(r), tg))
    end
  end

  -- Datos del motor y números grandes (con % no se puede escribir dentro de \directlua)
  function version_motor()
    tex.sprint(string.format("%s, motor %s %.2f", _VERSION, status.luatex_engine or "LuaTeX", status.luatex_version / 100))
  end
  function grandes()
    tex.sprint(string.format("$2^{64} = %.0f$ y $\\pi \\approx %.12f$", 2^64, math.pi))
  end

  -- Resistencias: serie, paralelo y divisor de tensión
  function circuito(vin, r1, r2)
    local par = r1 * r2 / (r1 + r2)
    tex.sprint(string.format(
      "Con $V_{in} = %g\\,\\mathrm{V}$, $R_1 = %g\\,\\Omega$ y $R_2 = %g\\,\\Omega$: "
      .. "serie $= %g\\,\\Omega$, paralelo $\\approx %.2f\\,\\Omega$ y "
      .. "$V_{out} = V_{in}\\,\\frac{R_2}{R_1 + R_2} \\approx %.3f\\,\\mathrm{V}$.",
      vin, r1, r2, r1 + r2, par, vin * r2 / (r1 + r2)))
  end
\end{luacode*}

\title{Prueba 1: cálculos con Lua}
\author{TeXsting}
\date{\today}

\begin{document}
\maketitle

Este PDF lo calcula Lua mientras compila: si ves números en vez de errores, el motor
LuaLaTeX funciona. Versión: \directlua{version_motor()}.

\section{Números}
Primos menores que 200: \directlua{primos(200)}.

\medskip
Primeros 25 de Fibonacci: \directlua{fibonacci(25)}.

\medskip
\directlua{grandes()}.

\section{Tabla trigonométrica generada}
\begin{center}
\begin{tabular}{rccc}
  \toprule
  Ángulo & $\sin$ & $\cos$ & $\tan$ \\
  \midrule
  \directlua{tabla_trig()}
  \bottomrule
\end{tabular}
\end{center}

\section{Un circuito}
\directlua{circuito(12, 1000, 2200)}

\section{Matemática con unicode-math}
Se escribe directo con símbolos Unicode:
\[
  ∫_0^∞ e^{-x^2}\,dx = \frac{\sqrt{π}}{2}, \qquad
  ∑_{n=1}^{∞} \frac{1}{n^2} = \frac{π^2}{6}, \qquad
  ∀ε>0\ ∃δ>0 : |x - a| < δ ⇒ |f(x) - f(a)| < ε .
\]
Y en negrita: $\symbf{F} = m\,\symbf{a}$, $∇ × \symbf{E} = -\frac{∂\symbf{B}}{∂t}$.

\section{Texto en español}
Acentos, eñes y signos: ¿Qué tal? ¡Muy bien! Pingüino, cigüeña, año, corazón.
\end{document}
